\section{Limitations}\label{sec:limitations}
\paragraph{Scope and finite-state assumption.}
This work addresses synthesis for a fixed target LLM agent, for which we assume that the event predicates are sufficient (A1). An evolving agent can break this assumption: new tools or behaviors can produce trajectories that the predicates map to the same word despite different policy judgments. The abstraction must then be refined and the contract re-learned, which we leave to future work on contract evolution. A high-level policy may also require unbounded identities, counters, data values, or obligations that no regular language over a practical finite alphabet captures (A2). We focus on prefix-closed safety properties, where a prohibited effect has an observable bad prefix and pre-effect interception is meaningful.

\paragraph{Assumptions and oracle fidelity.}
Our guarantee is conditional on (A1)--(A5). Policy-relevant observations can be missed, and predicates can merge trajectories with different ground-truth labels, violating (A1). An LLM with a pinned model and temperature zero is consistent by construction but may be stably incorrect, so (A3) can be measured but not enforced; its errors off the dataset enter the contract, and $L^*$ does not degrade gracefully when they do. Because labels win, label errors also enter the contract unless they collide with other labels. Without a complete dataset, the learned DFA is only guaranteed to agree with the dataset, and existing datasets are not complete in general. A dataset of trajectories also covers only realizable behavior, so even a complete one constrains only realizable words. Behavior absent from the dataset, such as a rare violating sequence, may be judged wrongly without any signal. Dataset labels (A5) can themselves be wrong, and a dataset collected from other agents may not reflect the target agent's behavior. Some abstract queries cannot be realized as actual agent trajectories, so the query bounds do not imply a bound on agent executions.

\paragraph{Enforcement boundary.}
The framework relies on agent-level instrumentation and assumes it can inspect relevant proposed effects before commitment. Missing adapters, incorrect metadata, concurrent actions, and partial tool failures may invalidate monitor state or allow an effect to occur before the trigger fires. A correct DFA cannot repair a missed enforcement point. We evaluate these practical failure modes and do not claim an OS-level or universal agent-safety guarantee.
